\chapter{Analytic manifolds}
\label{chap:analytic}

Real and complex analytic manifolds modeled on locally convex spaces, following
\cite[Definitions~3.1.3--3.1.4 and~2.2.18]{GN26}. Real analyticity is Milnor's notion: a map is
real analytic if it extends to a complex analytic map on an open subset of the complexification
$E_{\mathbb{C}} = E \times E$ (see the companion library \texttt{lc-analysis}). As for smooth manifolds,
no completeness of the model space is assumed: composition, locality, ``analytic implies
smooth'' and ``complex analytic implies real analytic'' hold for arbitrary Hausdorff locally
convex spaces.

\section{The analytic groupoids}

\begin{definition}[Analytic groupoids]
  \label{def:analytic-groupoids}
  \lean{LocallyConvex.Manifolds.realAnalyticGroupoid, LocallyConvex.Manifolds.complexAnalyticGroupoid,
    LocallyConvex.Manifolds.realAnalyticOn_congr}
  \leanok
  \uses{def:mb-groupoid}
  For a real locally convex space $E$, the \emph{real analytic groupoid} consists of the
  homeomorphisms between open subsets of $E$ which are real analytic, with real analytic
  inverse. For a complex locally convex space $E$, the \emph{complex analytic groupoid} is
  defined in the same way with complex analytic maps.
\end{definition}

\begin{proof}
  \leanok
  These are structure groupoids: compositions of real (complex) analytic maps are real
  (complex) analytic \cite[Proposition~2.2.4]{GN26}, and real (complex) analyticity is a local
  property \cite[Lemma~2.2.3]{GN26}.
\end{proof}

\begin{proposition}
  \label{prop:analytic-groupoids-le}
  \lean{LocallyConvex.Manifolds.realAnalyticGroupoid_le_mbGroupoid,
    LocallyConvex.Manifolds.complexAnalyticGroupoid_le_realAnalyticGroupoid}
  \leanok
  \uses{def:analytic-groupoids}
  Every real analytic diffeomorphism is a smooth diffeomorphism, and (for a complex locally
  convex space) every complex analytic diffeomorphism is real analytic.
\end{proposition}

\begin{proof}
  \leanok
  Real analytic maps are smooth \cite[Lemma~2.2.5]{GN26}; complex analytic maps are real
  analytic \cite[Proposition~2.2.11]{GN26}. Neither statement uses completeness.
\end{proof}

\begin{definition}[Analytic manifolds]
  \label{def:analytic-manifold}
  \lean{LocallyConvex.Manifolds.RealAnalyticManifold, LocallyConvex.Manifolds.ComplexAnalyticManifold,
    LocallyConvex.Manifolds.RealAnalyticManifold.mbManifold,
    LocallyConvex.Manifolds.ComplexAnalyticManifold.realAnalyticManifold}
  \leanok
  \uses{def:analytic-groupoids, prop:analytic-groupoids-le, def:mb-manifold}
  A \emph{real (complex) analytic manifold} is a charted space whose transition maps belong to
  the real (complex) analytic groupoid. By Proposition~\ref{prop:analytic-groupoids-le}, a
  complex analytic manifold is real analytic, and a real analytic manifold is smooth.
\end{definition}

\begin{remark}
  \label{rem:analytic-opens}
  Both groupoids are closed under restriction to open sets, so open subsets of analytic
  manifolds (in particular of the model space) are analytic manifolds
  \cite[Examples~3.1.8 and~3.1.9]{GN26}.
\end{remark}

\section{Analytic maps}

\begin{definition}[Analytic maps]
  \label{def:analytic-maps}
  \lean{LocallyConvex.Manifolds.RealAnalyticMOn, LocallyConvex.Manifolds.ComplexAnalyticMOn}
  \leanok
  \uses{def:analytic-manifold, def:contmdiff}
  A map $f \colon M \to N$ is \emph{real (complex) analytic} on an open set $U \subseteq M$ if
  it is continuous on $U$ and its local representatives in the preferred charts are real
  (complex) analytic.
\end{definition}

\begin{lemma}[Insertion of charts]
  \label{lem:chart-change}
  \lean{LocallyConvex.Manifolds.exists_chartChange}
  \leanok
  Near each point of the domain of $e' \circ f \circ e^{-1}$, this map factors as
  $(e' \circ c'^{-1}) \circ (c' \circ f \circ c^{-1}) \circ (c \circ e^{-1})$ with explicit
  open domains, for any other charts $c$, $c'$ \cite[Remark~3.1.5]{GN26}.
\end{lemma}

\begin{proof}
  \leanok
  Set theory.
\end{proof}

\begin{theorem}
  \label{thm:analytic-maps}
  \lean{LocallyConvex.Manifolds.realAnalyticOn_transition, LocallyConvex.Manifolds.RealAnalyticMOn.realAnalyticOn_chart,
    LocallyConvex.Manifolds.RealAnalyticMOn.comp, LocallyConvex.Manifolds.RealAnalyticMOn.smoothMBMOn,
    LocallyConvex.Manifolds.realAnalyticMOn_id, LocallyConvex.Manifolds.realAnalyticMOn_const,
    LocallyConvex.Manifolds.realAnalyticMOn_iff_realAnalyticOn,
    LocallyConvex.Manifolds.complexAnalyticOn_transition,
    LocallyConvex.Manifolds.ComplexAnalyticMOn.complexAnalyticOn_chart, LocallyConvex.Manifolds.ComplexAnalyticMOn.comp,
    LocallyConvex.Manifolds.ComplexAnalyticMOn.realAnalyticMOn, LocallyConvex.Manifolds.complexAnalyticMOn_id,
    LocallyConvex.Manifolds.complexAnalyticMOn_const, LocallyConvex.Manifolds.complexAnalyticMOn_iff_complexAnalyticOn}
  \leanok
  \uses{def:analytic-maps, lem:chart-change}
  The local representatives of a real (complex) analytic map in arbitrary charts of the
  maximal atlases are real (complex) analytic. Compositions of analytic maps are analytic;
  identities and constants are analytic; real analytic maps are smooth; complex analytic maps
  are real analytic; and between open subsets of the model spaces, the notions agree with those
  of \texttt{lc-analysis}.
\end{theorem}

\begin{proof}
  \leanok
  Insert charts (Lemma~\ref{lem:chart-change}) and use that the classes of analytic maps are
  closed under composition and local.
\end{proof}

\begin{example}[Inversion in a continuous inverse algebra]
  \label{ex:analytic-inversion}
  \lean{LocallyConvex.Manifolds.realAnalyticMOn_inverse}
  \leanok
  \uses{def:analytic-maps}
  The unit group $A^\times$ of a real continuous inverse algebra $A$ is open in $A$, hence a
  real analytic manifold, and inversion is real analytic on it \cite[Proposition~2.2.13]{GN26}.
\end{example}

\section{Complexifications}

\begin{definition}[Complexification]
  \label{def:complexification-manifold}
  \lean{LocallyConvex.Manifolds.IsComplexification, LocallyConvex.Manifolds.isComplexification_ofReal}
  \leanok
  \uses{def:analytic-manifold}
  A complex analytic manifold $M^*$ modeled on $E_{\mathbb{C}}$ is a \emph{complexification} of a real
  analytic manifold $M$ modeled on $E$ if $M \subseteq M^*$ and each $x \in M$ lies in the
  domain $U^*$ of a chart $\varphi \colon U^* \to U'$ of $M^*$ with
  $\varphi(M \cap U^*) = E \cap U'$ such that $\varphi|_{M \cap U^*}$ is a chart of $M$
  \cite[Definition~2.2.18]{GN26}. In Lean the inclusion is an injective map
  $\iota \colon M \to M^*$. The space $E_{\mathbb{C}}$ is a complexification of $E$.
\end{definition}

\begin{lemma}[Gluing near a closed set {\cite[Lemma~A.5.10]{GN26}}]
  \label{lem:paracompact-glue}
  \lean{LocallyConvex.Manifolds.exists_glue_of_paracompact}
  \leanok
  Let $Z$ be paracompact, $A \subseteq Z$ closed, $(V_i)$ open sets covering $A$, and $g_i$ maps
  such that $g_i$ and $g_j$ agree near each point of $A \cap V_i \cap V_j$. Then there are an open
  $U \supseteq A$ and a map $g$ which near each point of $U$ agrees with some $g_i$ on an open
  subset of $V_i$.
\end{lemma}

\begin{proof}
  \leanok
  As in \cite{GN26}: closed neighborhoods $Q_z \subseteq V_{i(z)}$ of the points of $A$, a locally
  finite open refinement $(W_j)$ of the cover by their interiors and $Z \setminus A$, and around
  each $a \in A$ a neighborhood meeting finitely many $W_j$, avoiding the closures of those not
  containing $a$ in their closure, and contained in the sets where the corresponding $g_i$ agree.
\end{proof}

\begin{theorem}[Extension to complexifications {\cite[Proposition~2.2.19]{GN26}}]
  \label{thm:complexification-extension}
  \lean{LocallyConvex.Manifolds.complexAnalyticMOn_of_mem_maximalAtlas,
    LocallyConvex.Manifolds.complexAnalyticMOn_symm_of_mem_maximalAtlas,
    LocallyConvex.Manifolds.complexAnalyticMOn_of_local,
    LocallyConvex.Manifolds.IsComplexification.exists_local_extension,
    LocallyConvex.Manifolds.IsComplexification.eventuallyEq_of_extension,
    LocallyConvex.Manifolds.IsComplexification.exists_complexAnalytic_extension}
  \leanok
  \uses{lem:paracompact-glue}
  Let $M^*$ be a complexification of the real analytic manifold $M$, with $M$ closed in $M^*$ and
  $M^*$ paracompact. Every real analytic map from $M$ to a complex analytic manifold $L$ extends
  to a complex analytic map on an open neighbourhood of $M$ in $M^*$.
\end{theorem}

\begin{proof}
  \leanok
  In a chart $\varphi$ of $M^*$ restricting to a chart $\psi$ of $M$ and a chart $\kappa$ of $L$,
  the local representative $\kappa \circ f \circ \psi^{-1}$ has a complex analytic extension, which
  after shrinking has real points only where $f$ is defined in these charts; this gives local
  extensions $g_x$. Two of them agree near each point of $M$: on a convex neighbourhood in
  $E_{\mathbb{C}}$ their representatives agree on the real points, hence everywhere by the identity
  theorem. Lemma~\ref{lem:paracompact-glue} glues them.
\end{proof}

\section{Real analytic forms}

\begin{definition}[Real analytic forms]
  \label{def:analytic-forms}
  \lean{LocallyConvex.Manifolds.DifferentialForm.IsRealAnalytic,
    LocallyConvex.Manifolds.DifferentialForm.analyticForms,
    LocallyConvex.Manifolds.SmoothVectorField.IsRealAnalytic}
  \leanok
  \uses{def:diff-form, def:analytic-manifold, def:interior}
  Let $M$ be a real analytic manifold. A form $\omega \in \Omega^n(M, F)$ is \emph{real
  analytic} if its local representatives (Definition~\ref{def:diff-form}) are real analytic in
  the sense of Milnor \cite[Remark~3.6.18]{GN26}. The real analytic forms are a subspace
  $\Omega^n_\omega(M, F)$ of $\Omega^n(M, F)$. A smooth vector field is real analytic if its local
  representatives are.
\end{definition}

\begin{theorem}[Operations on real analytic forms]
  \label{thm:analytic-forms}
  \lean{LocallyConvex.Manifolds.DifferentialForm.IsRealAnalytic.d,
    LocallyConvex.Manifolds.DifferentialForm.analyticDLinear,
    LocallyConvex.Manifolds.DifferentialForm.IsRealAnalytic.compCLM,
    LocallyConvex.Manifolds.DifferentialForm.IsRealAnalytic.wedge,
    LocallyConvex.Manifolds.DifferentialForm.IsRealAnalytic.pullback,
    LocallyConvex.Manifolds.SmoothVectorField.IsRealAnalytic.bracket,
    LocallyConvex.Manifolds.DifferentialForm.IsRealAnalytic.interior,
    LocallyConvex.Manifolds.DifferentialForm.IsRealAnalytic.lieDeriv}
  \leanok
  \uses{def:analytic-forms, thm:diff-form-d, def:diff-form-pullback, def:wedge, def:interior,
    def:lie-deriv, thm:analytic-maps}
  Exterior derivatives, compositions with continuous linear maps, wedge products and pullbacks
  along real analytic maps of real analytic forms are real analytic. Brackets of real analytic
  vector fields are real analytic, and so are $i_X\omega$ and $L_X\omega$ for real analytic $X$
  and $\omega$. Hence all the identities of Section~\ref{sec:lie-derivative} hold in
  $\Omega^\bullet_\omega(M, F)$, as asserted in \cite[Remark~3.6.18]{GN26}.
\end{theorem}

\begin{proof}
  \leanok
  In a chart, each operation is built from the local representatives by sums, compositions with
  continuous (multi)linear maps, compositions with real analytic maps and directional
  derivatives. Real analytic maps are closed under all of these, the last one because if $g$ is a
  complex analytic extension of $f$, then $\dd g$, read on $(E \times E)_{\mathbb{C}} \cong
  E_{\mathbb{C}} \times E_{\mathbb{C}}$, is a complex analytic extension of $\dd f$
  (\texttt{LocallyConvex.Analysis.RealAnalyticOn.dMB} in \texttt{lc-analysis}). For pullbacks,
  real analyticity is local \cite[Lemma~2.2.3]{GN26}, and near each point the representative of
  $\varphi^*\omega$ is a local pullback along a real analytic map between chart domains.
\end{proof}

\begin{definition}[Real analytic de Rham cohomology]
  \label{def:analytic-de-rham}
  \lean{LocallyConvex.Manifolds.DifferentialForm.analyticClosedForms,
    LocallyConvex.Manifolds.DifferentialForm.analyticExactForms,
    LocallyConvex.Manifolds.DifferentialForm.analyticExactForms_le_analyticClosedForms,
    LocallyConvex.Manifolds.AnalyticDeRham, LocallyConvex.Manifolds.AnalyticDeRham.toDeRham}
  \leanok
  \uses{thm:analytic-forms, def:de-rham}
  Let $Z^n_\omega(M, F)$ be the closed real analytic $n$-forms and $B^n_\omega(M, F) =
  \dd(\Omega^{n-1}_\omega(M, F))$, with $B^0_\omega = 0$. The real analytic de Rham cohomology is
  $H^n_\omega(M, F) = Z^n_\omega(M, F)/B^n_\omega(M, F)$, and there is a natural linear map
  $H^n_\omega(M, F) \to H^n_{dR}(M, F)$.
\end{definition}

\begin{remark}
  For a real analytic submanifold $V$ of $\R^n$ this map is an isomorphism
  \cite[\S7, consequence~(4), p.~89]{Car57}, a consequence of Cartan's Theorem~B for coherent
  analytic sheaves on analytic subsets of $\R^n$ \cite[Theorem~3]{Car57}; the proof is omitted
  there. This covers every second countable real analytic manifold, since such a manifold admits a
  real analytic, regular and proper embedding into some $\R^k$ \cite[Theorem~3, p.~470]{Gra58}. The proof uses
  complexifications and sheaf cohomology on Stein manifolds, with no counterpart for general
  locally convex models. We do not know whether it is injective or surjective in general.
\end{remark}

\begin{remark}[What is not covered]
  \label{rem:forms-not-covered}
  The following parts of \cite[Section~3.6]{GN26} are not formalized.
  \begin{itemize}
    \item The finite-dimensional material (Definition~3.6.19 to Remark~3.6.24: forms on subsets of
      $\R^n$, orientations, integration of forms and Stokes' Theorem).
  \end{itemize}
\end{remark}
